% ============================================================================
% Մանկավարժական լուծումներ - Դիսկրետ մաթեմատիկա, Գործնական 5
% (Կոմբինատորիկա, հաշվման հիմնական կանոնները). Խնդիրներ 1.3, 1.5, 1.10, 2.5, 2.10, 2.12.
% Հայերեն տարբերակ, XeLaTeX (fontspec + Sylfaen). Կազմել՝ xelatex-ով երկու անգամ։
% Տերմինաբանությունը՝ glossary_discrete_focused.md / glossary_Apr.csv-ի համաձայն։
% ============================================================================
\documentclass[11pt,a4paper]{article}

% ---- Տառատեսակ (XeLaTeX) ----
\usepackage{fontspec}
\setmainfont{Sylfaen}

% Հայերենը plain XeLaTeX-ում տողադարձման կանոններ չունի, ուստի երկար տողերը չեն
% կարող կոտրվել և դուրս կգան լուսանցքից։ Թող TeX-ը ձգի բառամիջյան բացատները։
\tolerance=2000
\emergencystretch=3.5em

% ---- Մաթեմատիկա ----
\usepackage{amsmath,amssymb,amsthm}
\usepackage{mathtools}

% ---- Դասավորություն ----
\usepackage[margin=2.3cm]{geometry}
\usepackage{enumitem}
\usepackage{booktabs}
\usepackage{array}

% ---- Գրաֆիկա ----
\usepackage{xcolor}
\usepackage[most]{tcolorbox}
\usepackage{tikz}
\usetikzlibrary{arrows.meta,positioning,calc,fit,backgrounds,shapes.misc}
\usepackage{pgfplots}
\pgfplotsset{compat=1.18}
\usepackage{hyperref}
\hypersetup{colorlinks=true,linkcolor=blue!55!black,urlcolor=blue!55!black}

% ---- Գունապնակ (Հայաստանի դրոշի գույները) ----
\definecolor{armRed}{HTML}{D90012}
\definecolor{armBlue}{HTML}{0033A0}
\definecolor{armOrange}{HTML}{F2A800}
\definecolor{ink}{HTML}{22303C}

\setlength{\parindent}{0pt}
\setlength{\parskip}{5pt}

% ---- Վանդակներ ----
\newtcolorbox{problembox}[1][]{enhanced, breakable, sharp corners=south,
  colback=armOrange!10, colframe=armOrange!75!black, boxrule=0.8pt, arc=2mm,
  fonttitle=\bfseries\color{white}, coltitle=white,
  attach boxed title to top left={xshift=4mm,yshift=-2.4mm},
  boxed title style={colback=armOrange!85!black, arc=1mm},
  title={Խնդիր #1}}

\newtcolorbox{ideabox}{enhanced, breakable,
  colback=armBlue!6, colframe=armBlue!70!black, boxrule=0.8pt, arc=2mm,
  fonttitle=\bfseries\color{white}, coltitle=white,
  attach boxed title to top left={xshift=4mm,yshift=-2.4mm},
  boxed title style={colback=armBlue!75!black, arc=1mm},
  title={Գաղափարը}}

\newtcolorbox{methodbox}{enhanced, breakable,
  colback=gray!8, colframe=gray!55!black, boxrule=0.7pt, arc=2mm,
  fonttitle=\bfseries, title={$\bigstar$\ Հիշարժան մեթոդ}}

\newtcolorbox{answerbox}{enhanced, sharp corners=south,
  colback=green!8, colframe=green!50!black, boxrule=1pt, arc=2mm, left=4mm,right=4mm}
\newcommand{\answer}[1]{\begin{answerbox}\centering\textbf{Պատասխան՝}\quad #1\end{answerbox}}

% ---- Վերնագիր ----
\title{\vspace{-1.2cm}\textbf{Դիսկրետ մաթեմատիկա - Գործնական 5}\\[2pt]
\large Կոմբինատորիկա, հաշվման հիմնական կանոնները՝ մանրամասն լուծումներ\\[2pt]
\normalsize\itshape դատողությունները՝ քայլ առ քայլ բացված}
\author{}
\date{}

\begin{document}
\maketitle
\vspace{-1.0cm}

\section*{Նախապատրաստում}
Այս թերթիկը \emph{առանց թվարկելու հաշվելու} մասին է։ Այստեղ գրեթե ամեն ինչ հենվում է երկու կանոնի վրա։
\textbf{Արտադրյալի կանոնը}՝ եթե ինչ-որ բան կառուցվում է առաջին ընտրությունը $a$ տարբերակով, ապա
անկախ կերպով երկրորդ ընտրությունը $b$ տարբերակով անելով, ապա ընդամենը կա $a\cdot b$ տարբերակ ($a$
շարք և $b$ սյունակ ունեցող ցանց)։ \textbf{Գումարման կանոնը}՝ եթե ինչ-որ բան բաժանվում է առանձին,
չհատվող դեպքերի, որոնք ունեն $a$ և $b$ տարբերակ, ապա ընդամենը կա $a+b$ տարբերակ։ Հետո երկու խնդիր
ավելացնում են \textbf{Դիրիխլեի սկզբունքը}՝ եթե $n+1$ առարկա դնես $n$ արկղում, ապա ինչ-որ արկղ
կպահի առնվազն երկուսը։ Հաշվումը տեսանելի է՝ նկարիր դիրքերը, դեպքերը կամ արկղերը, և բանաձևն ինքն իրեն
կկարդացվի նկարից։ Լուծում ենք 1 և 2 բաժինների 1.3, 1.5, 1.10, 2.5, 2.10, 2.12 համարակալված խնդիրները։

% ============================================================================
\section*{Խնդիր 1.3 - Տարբեր առարկաների երկու գիրք}

\begin{problembox}[1.3]
Դարակն ունի $15$ ինֆորմատիկայի գիրք, $12$ մաթեմատիկայի գիրք և $10$ ֆիզիկայի գիրք (բոլորը
տարբեր)։ Քանի՞ տարբերակով կարող ես ընտրել \textbf{երկու} \textbf{տարբեր} առարկայի գիրք։
\end{problembox}

\begin{ideabox}
«Տարբեր առարկաներ»-ը նշանակում է, որ երկու գրքերը գալիս են երկու տարբեր կույտից։ Այսպիսով՝ նախ որոշիր,
թե \emph{երեք առարկաներից որ երկուսն} է օգտագործում զույգը՝ կա այդպիսի երեք առարկա-զույգ՝
ինֆորմատիկա+մաթ, ինֆորմատիկա+ֆիզիկա, մաթ+ֆիզիկա։ Այս երեք դեպքերը երբեք չեն հատվում (զույգը
պատկանում է հենց դրանցից մեկին), ուստի \textbf{գումարման կանոնով} գումարում ենք դրանց քանակները։
Մեկ դեպքի ներսում, ասենք ինֆորմատիկա+մաթ, երկու ընտրությունները անկախ են՝ $15$ ինֆորմատիկայի
գրքերից յուրաքանչյուրը $12$ մաթ գրքերից յուրաքանչյուրի հետ, ինչը $15\cdot 12$ է \textbf{արտադրյալի
կանոնով}։ Այսպիսով՝ ամբողջ պատասխանը երեք արտադրյալների մեկ գումար է։
\end{ideabox}

\textbf{Քայլ 1 - Բաժանենք երեք առարկա-զույգերի (գումարման կանոն)։}
Վավեր զույգն օգտագործում է հենց դրանցից մեկը՝ $\{$ինֆ, մաթ$\}$, $\{$ինֆ, ֆիզ$\}$, $\{$մաթ, ֆիզ$\}$։

\textbf{Քայլ 2 - Հաշվենք յուրաքանչյուր զույգը (արտադրյալի կանոն)։}
\[
\underbrace{15\cdot 12}_{\text{ինֆ, մաթ}}=180,\qquad
\underbrace{15\cdot 10}_{\text{ինֆ, ֆիզ}}=150,\qquad
\underbrace{12\cdot 10}_{\text{մաթ, ֆիզ}}=120 .
\]

\textbf{Քայլ 3 - Գումարենք դեպքերը։}
\[
180+150+120 = 450 .
\]

\textbf{Ստուգում։} Զույգերը չկարգավորված են, բայց յուրաքանչյուր $15\cdot 12$ արտադրյալն արդեն հաշվում է
\emph{չկարգավորված} զույգը մեկ անգամ՝ «ինֆորմատիկայի գիրք և մաթ գիրք»-ը թվարկվում է մեկ անգամ, ոչ թե
երկու, քանի որ երկու գրքերը գալիս են տարբեր պիտակավորված կույտերից։ Կրկնակի հաշվում չկա, և երեք
դեպքերը չհատվող են։ \checkmark

\answer{$15\cdot12+15\cdot10+12\cdot10 = 180+150+120 = \mathbf{450}$}

\begin{center}
\begin{tikzpicture}[font=\small,
  pile/.style={draw, rounded corners=2pt, minimum width=2.5cm, minimum height=9mm, align=center},
  prod/.style={draw, rounded corners=2pt, minimum width=2.3cm, minimum height=8mm, align=center, fill=armBlue!10}]
% three piles
\node[pile, fill=armRed!12]    (I) at (0,0)   {Ինֆորմատիկա\\$15$};
\node[pile, fill=armBlue!12]   (M) at (3.3,0) {Մաթ\\$12$};
\node[pile, fill=armOrange!18] (P) at (6.6,0) {Ֆիզիկա\\$10$};
% the three subject-pairs and their products
\node[prod] (IM) at (0.4,-2.0)  {ինֆ $\times$ մաթ\\$15\cdot12=180$};
\node[prod] (IP) at (3.3,-2.0)  {ինֆ $\times$ ֆիզ\\$15\cdot10=150$};
\node[prod] (MP) at (6.2,-2.0)  {մաթ $\times$ ֆիզ\\$12\cdot10=120$};
\draw[armRed!70,thick]   (I) -- (IM); \draw[armBlue!70,thick]   (M) -- (IM);
\draw[armRed!70,thick]   (I) -- (IP); \draw[armOrange!80!black,thick] (P) -- (IP);
\draw[armBlue!70,thick]  (M) -- (MP); \draw[armOrange!80!black,thick] (P) -- (MP);
% sum
\node[draw, rounded corners=2pt, fill=green!12, minimum width=8.5cm, minimum height=8mm] (S) at (3.3,-3.6)
  {գումարման կանոն՝ $180+150+120=\mathbf{450}$};
\draw[->,>=Stealth] (IM) -- (S); \draw[->,>=Stealth] (IP) -- (S); \draw[->,>=Stealth] (MP) -- (S);
\end{tikzpicture}
\end{center}

\begin{methodbox}
\textbf{Փոխանցելի տեխնիկա՝} \emph{բազմապատկիր դեպքի ներսում, գումարիր դեպքերի վրայով։} Երբ առարկաները
բաժանվում են չհատվող տեսակների, յուրաքանչյուր տեսակը հաշվիր արտադրյալի կանոնով, ապա գումարիր
տեսակները (գումարման կանոն)։ Երեք խմբից երկուսն ընտրելը միշտ տալիս է հենց $\binom{3}{2}=3$ դեպք։
\end{methodbox}

% ============================================================================
\section*{Խնդիր 1.5 - 700-ից փոքր թվերը և բաժանելիությունը}

\begin{problembox}[1.5]
$700$-ից փոքր բնական թվերի մեջ՝ քանի՞սն են բաժանվում $5$-ի վրա։ Դրանցից քանի՞սն են բաժանվում նաև
$3$-ի վրա։ Քանի՞սն են բաժանվում և՛ $3$-ի, և՛ $5$-ի վրա։
\end{problembox}

\begin{ideabox}
Տրված սահմանից փոքր $d$-ի բազմապատիկները հավասարաչափ բաշխված են՝ $d,2d,3d,\dots$։ Ուստի «քանի՞ $d$-ի
բազմապատիկ կա, որ ամենաշատը $N$ է» պարզապես $\lfloor N/d\rfloor$ է՝ $N$-ի մեջ տեղավորում ես $d$ չափի
այնքան ամբողջ քայլ, որքան կարող ես։ «$700$-ից փոքր»-ը նշանակում է $1,2,\dots,699$, ուստի $N=699$։
Միակ նրբությունը վերջին հարցն է՝ «բաժանվում է և՛ $3$-ի, և՛ $5$-ի վրա»-ն նույնն է, ինչ «բաժանվում է
$\operatorname{lcm}(3,5)=15$-ի վրա», քանի որ $3$-ը և $5$-ը փոխադարձաբար պարզ են։ Իսկ «բաժանվում է
$5$-ի վրա և նաև $3$-ի վրա»-ն ճիշտ նույն պայմանն է՝ ուստի երկրորդ և երրորդ պատասխանները համընկնում են։
\end{ideabox}

\textbf{Քայլ 1 - Ֆիքսենք միջակայքը։}
$700$-ից փոքր բնական թվերն են $1,2,\dots,699$, ուստի օգտագործում ենք $N=699$։

\textbf{Քայլ 2 - Բաժանվող $5$-ի վրա։}
\[
\left\lfloor \tfrac{699}{5}\right\rfloor = 139 \qquad(\text{սրանք են } 5,10,15,\dots,695).
\]

\textbf{Քայլ 3 - Դրանցից նաև բաժանվող $3$-ի վրա։}
$5$-ի վրա և $3$-ի վրա բաժանվող թիվը բաժանվում է $\operatorname{lcm}(3,5)=15$-ի վրա (ճիշտ է, քանի որ
$\gcd(3,5)=1$)։ Ուստի հաշվում ենք $15$-ի բազմապատիկները՝
\[
\left\lfloor \tfrac{699}{15}\right\rfloor = 46 \qquad(\text{սրանք են } 15,30,\dots,690).
\]

\textbf{Քայլ 4 - Բաժանվող և՛ $3$-ի, և՛ $5$-ի վրա։}
Սա բառացիորեն նույն պայմանն է՝ բաժանվող $15$-ի վրա՝ ուստի քանակը կրկին $46$ է։

\textbf{Ստուգում։} $695 = 5\cdot 139 \le 699$ և հաջորդ բազմապատիկը՝ $700>699$, ուստի $139$-ը ճիշտ է։
$690 = 15\cdot 46 \le 699$ և $15\cdot 47 = 705 > 699$, ուստի $46$-ը ճիշտ է։ \checkmark

\answer{$\lfloor 699/5\rfloor = \mathbf{139}$-ը՝ բաժանվող $5$-ի վրա;\quad դրանցից $\lfloor 699/15\rfloor = \mathbf{46}$-ը՝ նաև $3$-ի վրա;\quad $\mathbf{46}$-ը՝ և՛ $3$-ի, և՛ $5$-ի վրա։}

\begin{center}
\begin{tikzpicture}[font=\small]
% number line with multiples of 5 ticked and multiples of 15 highlighted
\draw[->,>=Stealth] (0,0) -- (12.5,0) node[right]{};
\foreach \x/\lab in {0/0, 2/100, 4/200, 6/300, 8/400, 10/500, 11.9/695}{
  \draw (\x,0.12) -- (\x,-0.12);
}
\node[below] at (0,-0.15){$0$};
\node[below] at (6,-0.15){$\sim$};
\node[below] at (11.9,-0.18){$695$};
\node[below] at (12.5,0.05){$<700$};
% legend dots
\fill[armBlue] (1.0,0.55) circle (2.2pt); \node[right] at (1.15,0.55){$5$-ի բազմապատիկ (139 հատ)};
\fill[armRed]  (1.0,1.05) circle (2.8pt); \node[right] at (1.15,1.05){$15$-ի բազմապատիկ = $5$-ի և $3$-ի (46 հատ)};
% small inset showing every third multiple of 5 is a multiple of 15
\node[anchor=west] at (0,-1.1)
  {\footnotesize $5$-ի բազմապատիկների մեջ ամեն \emph{երրորդն} է նաև $3$-ի բազմապատիկ՝};
\foreach \i in {0,...,8}{
  \pgfmathsetmacro\xx{0.0+\i*1.35}
  \ifnum\i=0 \fill[armRed](\xx,-1.8)circle(3pt);
  \else\ifnum\i=3 \fill[armRed](\xx,-1.8)circle(3pt);
  \else\ifnum\i=6 \fill[armRed](\xx,-1.8)circle(3pt);
  \else \fill[armBlue](\xx,-1.8)circle(2pt); \fi\fi\fi
  \node[below,font=\scriptsize] at (\xx,-1.95){$\pgfmathparse{int(5+\i*5)}\pgfmathresult$};
}
\node[anchor=west,font=\scriptsize] at (0.0,-2.55){(կարմիր $=$ նաև $\div 3$՝ $15,30,45,\dots$)։ Հարաբերությունը՝ $139/3\approx 46$, համընկնում է $\lfloor 699/15\rfloor=46$-ի հետ։};
\end{tikzpicture}
\end{center}

\begin{methodbox}
\textbf{Փոխանցելի տեխնիկա՝} $\{1,\dots,N\}$-ում $d$-ի բազմապատիկների քանակը
$\lfloor N/d\rfloor$ է։ «Բաժանվող $a$-ի և $b$-ի վրա»-ն հանգում է
«բաժանվող $\operatorname{lcm}(a,b)$-ի վրա»-ի, իսկ փոխադարձաբար պարզ $a,b$-ի դեպքում դա
պարզապես $a\cdot b$ է։
\end{methodbox}

% ============================================================================
\section*{Խնդիր 1.10 - Յոթ թվանշանանի թվեր զույգ-կենտ ձևանմուշով}

\begin{problembox}[1.10]
Քանի՞ $7$-թվանշանանի թիվ ունի \textbf{կենտ} թվանշան յուրաքանչյուր \textbf{զույգ} դիրքում և
\textbf{զույգ} թվանշան յուրաքանչյուր \textbf{կենտ} դիրքում։ ($1$-ին դիրքն առաջատար թվանշանն է և չի
կարող լինել $0$.)
\end{problembox}

\begin{ideabox}
$7$-թվանշանանի թիվը պարզապես ձախից աջ լրացված յոթ դիրք է, և յուրաքանչյուր դիրքի թվանշանն ընտրվում է
մյուսներից \emph{անկախ}, ուստի սա մաքուր արտադրյալի-կանոնի հաշիվ է՝ բազմապատկում ենք յուրաքանչյուր
դիրքում թույլատրելի թվանշանների քանակը։ Կարդա կանոնը դիրք առ դիրք։ Կենտ դիրքերը՝ $1,3,5,7$, ուզում են
\emph{զույգ} թվանշան, որից կա հինգը՝ $\{0,2,4,6,8\}$։ Զույգ դիրքերը՝ $2,4,6$, ուզում են \emph{կենտ}
թվանշան, նույնպես հինգը՝ $\{1,3,5,7,9\}$։ Միակ խոչընդոտն ամենաառաջին դիրքն է՝ այն կենտ դիրք է, ուստի
ուզում է զույգ թվանշան, բայց չի կարող լինել $0$ (այլապես թիվն իրականում $7$-թվանշանանի չէ)։ Դա հանում է
հենց մեկ տարբերակ՝ թողնելով $\{2,4,6,8\}$, այսինքն՝ $4$ ընտրություն։ Մյուս բոլոր դիրքերը պահում են իրենց
ամբողջ $5$-ը։
\end{ideabox}

\textbf{Քայլ 1 - Դասակարգենք դիրքերը։}
$1$-ից $7$ դիրքերը։ \emph{Կենտ} դիրքերը՝ $\{1,3,5,7\}$, վերցնում են զույգ թվանշան, իսկ \emph{զույգ}
դիրքերը՝ $\{2,4,6\}$, վերցնում են կենտ թվանշան։

\textbf{Քայլ 2 - Հաշվենք յուրաքանչյուր դիրքի ընտրությունները։}
Զույգ թվանշաններ $\{0,2,4,6,8\}$՝ $5$ ամեն մեկում։ Կենտ թվանշաններ $\{1,3,5,7,9\}$՝ $5$ ամեն մեկում։
$1$-ին դիրքը հատուկ է՝ զույգ թվանշան, բայց ոչ $0$, ուստի $\{2,4,6,8\}$, ընդամենը $4$ ընտրություն։

\textbf{Քայլ 3 - Բազմապատկենք (արտադրյալի կանոն)։}
$4$ ընտրությամբ մեկ դիրք ($1$-ին դիրքը) և մյուս վեց դիրքերը՝ $5$-ական ընտրությամբ՝
\[
\underbrace{4}_{\text{դիրք }1}\;\cdot\;\underbrace{5\cdot5\cdot5}_{\text{դիրք }3,5,7}\;\cdot\;\underbrace{5\cdot5\cdot5}_{\text{դիրք }2,4,6}
= 4\cdot 5^{6} = 4\cdot 15625 = 62500 .
\]

\textbf{Ստուգում։} Եթե $1$-ին դիրքը կարողանար լինել նաև $0$, կստանայինք $5^7=78125$՝ առաջատար $0$-ի
արգելքը հանում է հենց այն թվերը, որոնք $1$-ին դիրքում ունեն $0$, ինչը $\tfrac15$ բաժին է, թողնելով
$\tfrac45\cdot 5^7 = 4\cdot 5^6 = 62500$։ Նույն թիվը՝ ստացված երկրորդ ճանապարհով։ \checkmark

\answer{$4\cdot 5^{6} = \mathbf{62500}$}

\begin{center}
\begin{tikzpicture}[font=\small,
  slot/.style={draw, minimum width=1.55cm, minimum height=1.15cm, align=center},
  odd/.style={slot, fill=armBlue!10},   % odd position -> even digit
  evn/.style={slot, fill=armOrange!16}, % even position -> odd digit
  first/.style={slot, fill=armRed!14, very thick}]
\node[first] (p1) at (0,0)   {զույգ,\\ոչ $0$\\\textbf{4}};
\node[evn]   (p2) at (1.6,0) {կենտ\\\textbf{5}};
\node[odd]   (p3) at (3.2,0) {զույգ\\\textbf{5}};
\node[evn]   (p4) at (4.8,0) {կենտ\\\textbf{5}};
\node[odd]   (p5) at (6.4,0) {զույգ\\\textbf{5}};
\node[evn]   (p6) at (8.0,0) {կենտ\\\textbf{5}};
\node[odd]   (p7) at (9.6,0) {զույգ\\\textbf{5}};
\foreach \i/\x in {1/0,2/1.6,3/3.2,4/4.8,5/6.4,6/8.0,7/9.6}{
  \node[above,font=\scriptsize\bfseries] at (\x,0.62){դիրք \i};
}
% digit-set reminders
\node[below,font=\scriptsize,align=center] at (0,-0.72){$\{2,4,6,8\}$};
\node[below,font=\scriptsize,align=center] at (1.6,-0.72){$\{1,3,5,7,9\}$};
\node[below,font=\scriptsize,align=center] at (3.2,-0.72){$\{0,2,4,6,8\}$};
\node[below,font=\scriptsize,align=center] at (4.8,-0.72){$\{1,3,5,7,9\}$};
\node[below,font=\scriptsize,align=center] at (6.4,-0.72){$\{0,2,4,6,8\}$};
\node[below,font=\scriptsize,align=center] at (8.0,-0.72){$\{1,3,5,7,9\}$};
\node[below,font=\scriptsize,align=center] at (9.6,-0.72){$\{0,2,4,6,8\}$};
% product line
\node[anchor=center,font=\small] at (4.8,-1.85)
  {$4\times 5\times 5\times 5\times 5\times 5\times 5 \;=\; 4\cdot 5^{6} \;=\; \mathbf{62500}$};
\node[anchor=center,font=\scriptsize\itshape] at (4.8,-2.45)
  {կապույտ $=$ կենտ դիրք (զույգ թվանշան), նարնջագույն $=$ զույգ դիրք (կենտ թվանշան), կարմիր $=$ առաջատար դիրք};
\end{tikzpicture}
\end{center}

\begin{methodbox}
\textbf{Փոխանցելի տեխնիկա՝} «լրացրու դիրքերը» տիպի խնդիրների համար, որտեղ
յուրաքանչյուր դիրք ընտրվում է անկախ, բազմապատկիր ամեն դիրքի ընտրությունների քանակները։ Առաջատար
թվանշանի «ոչ $0$» սահմանափակումը մշակիր այդ մեկ դիրքի քանակը նվազեցնելով։
\end{methodbox}

% ============================================================================
\section*{Խնդիր 2.5 - Վեց տղա և յոթ աղջիկ շարքում, ծայրերը՝ տղաներ}

\begin{problembox}[2.5]
Վեց տղա և յոթ աղջիկ ($13$ տարբեր մարդ) կանգնում են $13$ տեղանի շարքում։ Քանի՞ տարբերակով կարող են
շարվել այնպես, որ \textbf{առաջին և վերջին} տեղերը երկուսն էլ զբաղեցնեն տղաներ։
\end{problembox}

\begin{ideabox}
Կարգը կարևոր է (շարք), ուստի սա դասավորությունների հաշիվ է։ Մաքուր քայլն այն է, որ \emph{նախ
տեղադրես սահմանափակված դիրքերը}՝ ամեն ինչից առաջ տեղադրիր երկու ծայրային դիրքերը, որոնք հարկադրաբար
տղաներ են։ Ընտրիր ձախ ծայրի տղային ($6$ ընտրություն), ապա աջ ծայրի տղային ($5$ ընտրություն մնաց),
ուստի՝ $6\cdot 5$ տարբերակ երկու ծայրերի համար։ Հիմա երկու ծայրերը կարգավորված են, և սահմանափակումը
լիովին օգտագործված է՝ մնացած $11$ մարդը (մյուս $4$ տղան և բոլոր $7$ աղջիկները) լրիվ ազատ են լրացնելու
միջին $11$ տեղերը ցանկացած կարգով, ինչը $11!$ տարբերակ է։ Բազմապատկում ենք անկախ փուլերը։
\end{ideabox}

\textbf{Քայլ 1 - Լրացնենք ծայրերը (միակ սահմանափակումը)։}
Ձախ ծայր՝ $6$ տղաներից յուրաքանչյուրը։ Աջ ծայր՝ մնացած $5$ տղաներից յուրաքանչյուրը։ Դա $6\cdot 5 = 30$
կարգավորված տարբերակ է՝ սա $A_6^2 = 6\cdot 5$ դասավորությունն է։

\textbf{Քայլ 2 - Լրացնենք միջինն ազատ։}
Ծայրերը ֆիքսելուց հետո $11$ տարբեր մարդ է մնում $11$ միջին տեղերի համար՝ $11!$ կարգավորում։

\textbf{Քայլ 3 - Բազմապատկենք (արտադրյալի կանոն)։}
\[
6\cdot 5\cdot 11! = 30\cdot 39\,916\,800 = 1\,197\,504\,000 .
\]

\textbf{Ստուգում։} Առանց սահմանափակման ընդհանուր դասավորությունները կլինեին $13!$։ Մեր պայմանը
հարկադրում է, որ երկու ծայրերը լինեն $6$ տղաների մեջ՝ բոլոր $13!$ շարքերի այն բաժինը, որ ամեն ծայրում
ունի տղա, $\tfrac{6}{13}\cdot\tfrac{5}{12}$ է, ինչը տալիս է $13!\cdot\tfrac{6\cdot5}{13\cdot12} = \tfrac{13!}{156}\cdot30$։
Քանի որ $13! = 6\,227\,020\,800$, սա $6\,227\,020\,800\cdot\tfrac{30}{156}=1\,197\,504\,000$ է։ Նույն թիվը։ \checkmark

\answer{$6\cdot 5\cdot 11! = 30\cdot 11! = \mathbf{1\,197\,504\,000}$}

\begin{center}
\begin{tikzpicture}[font=\small,
  box/.style={draw, minimum width=0.78cm, minimum height=0.95cm},
  endb/.style={box, fill=armRed!16, very thick},
  mid/.style={box, fill=armBlue!8}]
\node[endb] (e1) at (0,0) {Տ};
\foreach \i in {1,...,11}{
  \node[mid] (m\i) at (\i*0.85,0) {};
}
\node[endb] (e2) at (12*0.85,0) {Տ};
% labels
\node[above,font=\scriptsize\bfseries] at (0,0.6){դիրք 1};
\node[above,font=\scriptsize\bfseries] at (12*0.85,0.6){դիրք 13};
\node[above,font=\scriptsize] at (6*0.85,0.6){դիրքեր $2$-ից $12$};
% braces
\draw[decorate,decoration={brace,amplitude=4pt},armRed!80!black]
  (-0.42,-0.62) -- (0.42,-0.62) node[midway,below,font=\scriptsize,yshift=-2pt]{$6$};
\draw[decorate,decoration={brace,amplitude=4pt},armRed!80!black]
  (12*0.85-0.42,-0.62) -- (12*0.85+0.42,-0.62) node[midway,below,font=\scriptsize,yshift=-2pt]{$5$};
\draw[decorate,decoration={brace,amplitude=5pt},armBlue!80!black]
  (0.85-0.42,-0.62) -- (11*0.85+0.42,-0.62) node[midway,below,font=\scriptsize,yshift=-3pt]{$11!$ (մյուս 11 մարդը)};
\node[font=\small] at (6*0.85,-1.85){$\underbrace{6\cdot 5}_{\text{երկու տղա-ծայրեր}}\;\cdot\;\underbrace{11!}_{\text{միջին}}=\mathbf{1\,197\,504\,000}$};
\end{tikzpicture}
\end{center}

\begin{methodbox}
\textbf{Փոխանցելի տեխնիկա՝} \emph{նախ տեղադրիր սահմանափակված դիրքերը}, ապա մնացածը դասավորիր ազատ։
$n$-ից $k$ անվանված դիրքերի համար կարգավորված մարդիկ ընտրելը $A_n^k = n(n-1)\cdots(n-k+1)$ է՝
մնացած $n-k$ մարդը մնացած դիրքերը լրացնում է $(n-k)!$ տարբերակով։
\end{methodbox}

% ============================================================================
\section*{Խնդիր 2.10 - Երեք փոխադարձ ընկերներ կամ երեք փոխադարձ անծանոթներ}

\begin{problembox}[2.10]
$6$ մարդանոց սենյակում նրանցից յուրաքանչյուր երկուսը կա՛մ \textbf{ընկերներ} են, կա՛մ
\textbf{անծանոթներ}։ Ապացուցիր, որ միշտ կա երեք մարդ, որոնք զույգ առ զույգ ընկերներ են, կամ երեք
մարդ, որոնք զույգ առ զույգ անծանոթներ են։
\end{problembox}

\begin{ideabox}
Մոդելավորիր այն որպես լրիվ գրաֆ $K_6$՝ $6$ կետ, յուրաքանչյուր զույգի միջև մեկ կող, յուրաքանչյուր կող
ներկված \textbf{կարմիր} (ընկերներ) կամ \textbf{կապույտ} (անծանոթներ)։ Պնդումն այն է, որ ինչ-որ
եռանկյուն միագույն է։ Հնարքն այն է, որ նայենք \emph{մեկ} մարդ $P$-ից։ $P$-ն ունի $5$ կող, երկու գույն՝
\textbf{Դիրիխլեի սկզբունքով} դրանցից առնվազն $\lceil 5/2\rceil = 3$-ն ունեն նույն գույնը՝ ասենք $P$-ն
ընկեր է (կարմիր) երեք մարդկանց՝ $A,B,C$-ի հետ։ Հիմա ամեն ինչ կախված է $A,B,C$ եռանկյունից։
Եթե $AB,BC,CA$ զույգերից թեկուզ մեկը կարմիր է, ապա այդ զույգը գումարած $P$-ն ամբողջովին կարմիր
եռանկյուն է (երեք փոխադարձ ընկերներ)։ Եթե դրանցից \emph{ոչ մեկը} կարմիր չէ, ապա $AB,BC,CA$-ն բոլորը
կապույտ են, ուստի $A,B,C$-ն իրենք ամբողջովին կապույտ եռանկյուն են (երեք փոխադարձ անծանոթներ)։
Փախուստ չկա։ Սա հենց $R(3,3)\le 6$ պնդումն է։
\end{ideabox}

\textbf{Քայլ 1 - Կառուցենք մոդելը։}
Գագաթներ $=$ $6$ մարդիկ՝ ներկիր $\binom{6}{2}=15$ կողերից յուրաքանչյուրը կարմիր (ընկերներ) կամ
կապույտ (անծանոթներ)։ «Երեք փոխադարձ ընկերներ» $=$ կարմիր եռանկյուն, «երեք
փոխադարձ անծանոթներ» $=$ կապույտ եռանկյուն։

\textbf{Քայլ 2 - Դիրիխլե մեկ գագաթում։}
Ֆիքսիր ցանկացած մարդ $P$։ $P$-ից դուրս եկող $5$ կողերը ներկված են $2$ գույնով, և
$\lceil 5/2\rceil=3$, ուստի Դիրիխլեի սկզբունքով դրանցից առնվազն $3$-ն ունեն նույն գույնը։ Առանց
ընդհանրությունը խախտելու՝ $P$-ն \textbf{կարմիրով} միացած է երեք մարդկանց՝ $A,B,C$-ի (եթե փոխարենը
մեծամասնության գույնը կապույտ է, ամենուր փոխիր «ընկերներ»/«անծանոթներ» բառերը՝
դատողությունը նույնն է)։

\textbf{Քայլ 3 - Քննենք $A,B,C$ եռանկյունը (երկու դեպք)։}
\begin{itemize}[leftmargin=1.3em,topsep=2pt]
\item \emph{$A,B,C$-ի մեջ ինչ-որ կող կարմիր է}, ասենք $AB$։ Այդ դեպքում $P\!A$, $P\!B$, $AB$-ն բոլորը կարմիր են՝
$P,A,B$-ն երեք փոխադարձ ընկերներ են։ Ավարտ։
\item \emph{$A,B,C$-ի մեջ ոչ մի կող կարմիր չէ։} Այդ դեպքում $AB$, $BC$, $CA$-ն բոլորը կապույտ են, ուստի $A,B,C$-ն
երեք փոխադարձ անծանոթներ են։ Ավարտ։
\end{itemize}
Երկու դեպքերից մեկը պարտադիր տեղի է ունենում, ուստի միագույն եռանկյունը միշտ գոյություն ունի։ $\blacksquare$

\textbf{Ստուգում (արդյո՞ք $6$-ը ճիշտ թիվն է)։} Միայն $5$ մարդով պնդումը կարող է ձախողվել՝ ներկիր
$K_5$-ը որպես կարմիր $5$-ցիկլ և կապույտ $5$-ցիկլ («հնգանկյունն ու
հնգաստղ») և ոչ մի եռանկյուն միագույն չէ։ Ուստի $5$-ը բավարար չէ, և $6$-ը սուր է՝
ահա թե ինչու է խնդիրն ասում $6$։ (Ուղիղ ճշգրիտ ստուգումը հաստատում է՝ $K_6$-ի բոլոր $2^{15}$ ներկումների մեջ
\emph{զրոն} է խուսափում միագույն եռանկյունից, մինչդեռ $K_5$-ն ունի ներկումներ, որոնք խուսափում են։) \checkmark

\answer{Այդպիսի միագույն եռանկյունը միշտ գոյություն ունի՝ սա $R(3,3)\le 6$ է։}

\begin{center}
\begin{tikzpicture}[font=\small, scale=1.0,
  ppl/.style={circle, draw, fill=white, inner sep=1.3pt, minimum size=6mm, font=\small\bfseries}]
% Left picture: P with 5 edges, 3 red by pigeonhole
\begin{scope}[xshift=-3.2cm]
\node[ppl, fill=armOrange!25] (P) at (0,0) {$P$};
\node[ppl] (A) at (90:1.9)  {$A$};
\node[ppl] (B) at (162:1.9) {$B$};
\node[ppl] (C) at (234:1.9) {$C$};
\node[ppl] (D) at (306:1.9) {$D$};
\node[ppl] (E) at (18:1.9)  {$E$};
% 3 red edges P-A,P-B,P-C ; 2 other edges blue (P-D, P-E)
\draw[armRed,very thick] (P)--(A);
\draw[armRed,very thick] (P)--(B);
\draw[armRed,very thick] (P)--(C);
\draw[armBlue,thick,dashed] (P)--(D);
\draw[armBlue,thick,dashed] (P)--(E);
\node[font=\scriptsize, align=center] at (0,-2.7)
  {$5$ կող, $2$ գույն $\Rightarrow$\\ $\ge 3$-ն ունեն նույն գույնը (այստեղ կարմիր՝ $A,B,C$)};
\end{scope}
% Right picture: triangle ABC case split
\begin{scope}[xshift=3.4cm]
\node[ppl, fill=armOrange!25] (P2) at (0,1.5) {$P$};
\node[ppl] (A2) at (-1.5,-0.6) {$A$};
\node[ppl] (B2) at (1.5,-0.6)  {$B$};
\node[ppl] (C2) at (0,-1.9)    {$C$};
\draw[armRed,very thick] (P2)--(A2);
\draw[armRed,very thick] (P2)--(B2);
\draw[armRed,very thick] (P2)--(C2);
% the three ABC edges: if any is red -> red triangle with P; else all blue -> blue triangle ABC
\draw[armRed,very thick] (A2)--(B2) node[midway,above,font=\scriptsize,yshift=1pt]{կարմի՞ր};
\draw[armBlue,thick,dashed] (B2)--(C2);
\draw[armBlue,thick,dashed] (C2)--(A2);
\node[font=\scriptsize, align=center] at (0,-3.0)
  {եթե $AB,BC,CA$-ից որևէ մեկը կարմիր է՝ կարմիր եռանկյուն $P$-ի հետ,\\ եթե բոլոր երեքը կապույտ են՝ $A,B,C$ ամբողջովին կապույտ եռանկյուն};
\end{scope}
% legend
\node[font=\scriptsize, anchor=west] at (-1.3,-4.1) {\textcolor{armRed}{\rule{8mm}{1.1pt}} կարմիր = ընկերներ};
\node[font=\scriptsize, anchor=west] at (1.7,-4.1) {\textcolor{armBlue}{\rule{8mm}{1.1pt}} կապույտ (կետագիծ) = անծանոթներ};
\end{tikzpicture}
\end{center}

\begin{methodbox}
\textbf{Փոխանցելի տեխնիկա (Ռամսեյ Դիրիխլեի միջոցով)՝} կառուցվածք հարկադրելու համար կենտրոնացիր մեկ
գագաթի վրա, Դիրիխլեով դրա կողերը բաշխիր մեծամասնության գույնի մեջ, ապա ստուգիր այն փոքր խումբը, որ
ստեղծում է մեծամասնությունը։ Երկու դեպքերի բաժանումը՝ «ներսում ինչ-որ կող $X$ գույնի է»-ն ընդդեմ
«ոչ մեկը»-ի, շարժիչն է։ Սա ապացուցում է $R(3,3)\le 6$։
\end{methodbox}

% ============================================================================
\section*{Խնդիր 2.12 - $7,77,777,\dots$-ից մեկը բաժանվում է $2013$-ի վրա}

\begin{problembox}[2.12]
Դիտարկիր $7,\,77,\,777,\,7777,\dots$ հաջորդականությունը ($n$-րդ անդամը $n$ յոթերով գրված թիվն է)։
Ապացուցիր, որ նրա առաջին $2014$ անդամների մեջ առնվազն մեկը բաժանվում է $2013$-ի վրա։
\end{problembox}

\begin{ideabox}
Չես կարող մատնանշել, թե որ անդամն է աշխատում առանց հաշվելու, բայց դրա կարիքը չունես՝ պարզապես պետք է
ապացուցես, որ մեկը \emph{գոյություն ունի}։ Գործիքը \textbf{Դիրիխլեի սկզբունքն} է՝ կիրառված մնացորդների
վրա։ Նայիր առաջին $2014$ անդամները $2013$ մոդուլով։ $2013$ մոդուլով մնացորդը $\{0,1,\dots,2012\}$
$2013$ արժեքներից մեկն է՝ դա $2013$ «արկղ» է։ Բայց մենք ունենք $2014$ անդամ դրանց մեջ գցելու
համար, ուստի երկու անդամ պետք է ընկնեն նույն արկղում, այսինքն՝ ունենան նույն մնացորդը։ Դրանց
\emph{տարբերությունն} այդ դեպքում բաժանվում է $2013$-ի վրա։ Վերջին քայլն այն տեսնելն է, որ այս
տարբերությունն ինքն իսկ ճիշտ ձևի թիվ է (յոթեր, որոնց հաջորդում են զրոներ), և զրոները ($10$-ի աստիճան)
կարելի է հեռացնել, քանի որ $2013$-ը փոխադարձաբար պարզ է $10$-ի հետ։ Մնում է հաջորդականության այն
անդամը, որ բաժանվում է $2013$-ի վրա։
\end{ideabox}

\textbf{Քայլ 1 - Կառուցենք արկղերն ու առարկաները։}
Թող $a_n = \underbrace{77\cdots7}_{n\text{ յոթ}}$։ Նայիր մնացորդները՝
$a_1 \bmod 2013,\ a_2\bmod 2013,\ \dots,\ a_{2014}\bmod 2013$։ Յուրաքանչյուր մնացորդ $0,1,\dots,2012$
$2013$ արժեքներից մեկն է։

\textbf{Քայլ 2 - Դիրիխլե։}
Ունենք $2014$ մնացորդ, բայց ընդամենը $2013$ հնարավոր արժեք, ուստի դրանցից երկուսը հավասար են՝
\[
a_i \equiv a_j \pmod{2013},\qquad \text{ինչ-որ } 1\le i < j \le 2014 \text{-ի համար} .
\]
(Եթե փոխարենը մնացորդներից մեկն արդեն $0$ է, ապա այդ անդամը բաժանվում է $2013$-ի վրա, և մենք
անմիջապես ավարտում ենք։)

\textbf{Քայլ 3 - Տարբերությունը բաժանվում է $2013$-ի վրա։}
$a_i\equiv a_j$-ից ստանում ենք $2013 \mid (a_j - a_i)$։ Հանելով յոթ-միասեռ թվերը՝
\[
a_j - a_i = \underbrace{77\cdots7}_{j} - \underbrace{77\cdots7}_{i}
= \underbrace{77\cdots7}_{\,j-i\,}\underbrace{00\cdots0}_{\,i\,}
= a_{\,j-i}\cdot 10^{\,i}.
\]
(Վերին $j-i$ յոթերը գոյատևում են, ստորին $i$ թվանշանները չեղարկվում են զրոների։)

\textbf{Քայլ 4 - Հեռացնենք $10$-ի աստիճանը։}
$2013 = 3\cdot 11\cdot 61$, որը $10=2\cdot 5$-ի հետ ընդհանուր արտադրիչ չունի, ուստի $\gcd(2013,10^{\,i})=1$։
Քանի որ $2013 \mid a_{\,j-i}\cdot 10^{\,i}$ և $2013$-ը փոխադարձաբար պարզ է $10^{\,i}$-ի հետ, այն պետք է
բաժանի մյուս արտադրիչը՝
\[
2013 \mid a_{\,j-i}.
\]
Վերջապես $1 \le j-i \le 2013 < 2014$, ուստի $a_{\,j-i}$-ն առաջին $2014$ անդամներից մեկն է։ $\blacksquare$

\textbf{Ստուգում։} Ապացույցը երաշխավորում է միայն \emph{գոյությունը}՝ ուղիղ հաշվարկը ցույց է տալիս, որ
իրականում աշխատող առաջին անդամը $a_{60}$-ն է (վաթսուն $7$-ով թիվը), և իսկապես $60 \le 2014$։
Ուստի գոյության պնդումը ոչ միայն ճշմարիտ է, այլև հարմարավետ կերպով։ \checkmark

\answer{Ինչ-որ $a_{j-i}$ $1\le j-i\le 2013$-ով բաժանվում է $2013$-ի վրա՝ իրականում $a_{60}$-ն արդեն բաժանվում է։}

\begin{center}
\begin{tikzpicture}[font=\small,
  trm/.style={draw, rounded corners=2pt, minimum width=1.05cm, minimum height=7mm, fill=armBlue!10, font=\scriptsize},
  box/.style={draw, minimum width=1.15cm, minimum height=9mm, fill=armOrange!12}]
% top row: 2014 terms (show a few then dots)
\node[font=\scriptsize] at (-1.3,1.6) {$2014$ անդամ՝};
\node[trm] (t1) at (0,1.6) {$a_1$};
\node[trm] (t2) at (1.1,1.6) {$a_2$};
\node[trm] (t3) at (2.2,1.6) {$a_3$};
\node at (3.0,1.6) {$\cdots$};
\node[trm] (tk) at (4.0,1.6) {$a_i$};
\node at (4.85,1.6) {$\cdots$};
\node[trm] (tl) at (5.8,1.6) {$a_j$};
\node at (6.65,1.6) {$\cdots$};
\node[trm] (tn) at (7.7,1.6) {$a_{2014}$};
% boxes: 2013 remainder boxes
\node[font=\scriptsize] at (-1.3,-0.3) {$2013$ արկղ՝};
\node[box] (b0) at (0,-0.3) {$0$};
\node[box] (b1) at (1.2,-0.3) {$1$};
\node[box] (b2) at (2.4,-0.3) {$2$};
\node at (3.4,-0.3) {$\cdots$};
\node[box] (br) at (4.6,-0.3) {$r$};
\node at (5.6,-0.3) {$\cdots$};
\node[box] (be) at (7.0,-0.3) {$2012$};
% two terms a_i, a_j land in same box r
\draw[->,>=Stealth,armRed,thick] (tk) to[bend right=10] (br.north);
\draw[->,>=Stealth,armRed,thick] (tl) to[bend left=18] (br.north);
\node[font=\scriptsize,armRed!80!black,align=center] at (4.6,-1.4)
  {երկու անդամ, նույն մնացորդը՝ $r$};
\node[font=\scriptsize,align=center] at (4.6,-2.15)
  {$\Rightarrow 2013\mid(a_j-a_i)=a_{j-i}\cdot 10^{\,i}$, և $\gcd(2013,10)=1\Rightarrow 2013\mid a_{j-i}$};
\node[font=\scriptsize] at (3.5,2.35) {$2014$ աղավնի $2013$ բնում $\Rightarrow$ բախում};
\end{tikzpicture}
\end{center}

\begin{methodbox}
\textbf{Փոխանցելի տեխնիկա՝} «ինչ-որ անդամ բաժանվում է $m$-ի վրա» ցույց տալու համար
վերցրու $m+1$ թեկնածու և Դիրիխլեով դասավորիր դրանց մնացորդները $m$ մոդուլով՝ կրկնությունը տալիս է
տարբերություն, որ բաժանվում է $m$-ի վրա։ Եթե տարբերությունը կրում է $10^{k}$ արտադրիչ, հեռացրու այն՝
օգտագործելով $\gcd(m,10)=1$։
\end{methodbox}

\end{document}
